% Direction A, full text, v2, 8pt-minimum variant (2026-11-02). Source: mock_A_full_v2.tex (unchanged).
% Every piece of text is \footnotesize (8pt) or larger; icon glyphs are drawn, not typeset (mock_macros_8pt.tex).
% Double-wide (figure*, 516pt) "snake" layout sharing the left sources-of-truth lane:
%   row 1: Stage 1  --> Stage 1.5        (output 1 crosses the gap between them)
%   row 2: Stage 3 <--  Stage 2          (output 2 drops from 1.5 into 2; output 3 sits under Stage 2
%                                         and feeds Stage 3 through the gutter, entering it from the top)
% Stage 3 sits next to the lane, so its two sources keep short arrows and the dashed SimSpect-goal
% arrow stays a vertical line in the lane, as in the one-column version.
% Changes from v2 besides font size: Stage 1.5's "+4. Pin" slab moved from the top of the list to the
% bottom (just above the solver), so output 1 can enter it from the side; the Stage-1 collector bar
% enters the solver from above (the solver's right end is now the exit).
\documentclass[tikz,border=0pt]{standalone}
\usepackage{newtxtext,newtxmath}
\usepackage{graphicx}
% 8pt-minimum variant of mock_macros.tex (2026-11-02): every text style is \footnotesize (8pt);
% icon glyphs that were \tiny text ({ } and </>) are drawn as paths instead.
% Shared shape vocabulary for the big-figure design mocks (BIGFIG_SHAPE_DESIGN.md).
% Every tikzpicture that uses these must set [x=1pt,y=1pt]: all numbers below are pt.
\usetikzlibrary{calc,fit,shapes.symbols,shapes.misc,shapes.geometric,decorations.pathmorphing,arrows.meta}
\providecolor{sssrcfill}{gray}{0.88}\providecolor{sssrcbg}{gray}{0.88}\providecolor{sssrcink}{gray}{0}\providecolor{ssxmitink}{gray}{0}
\providecolor{ssspecfill}{gray}{0.95}\providecolor{ssprotfill}{gray}{0.86}\providecolor{ssxmitfill}{gray}{0.9}
\def\defshade{sssrcbg}   % the one color: defense-specific / sources of truth (palette ../sscolors.tex)
\tikzset{
  >=latex,
  every node/.style={font=\footnotesize, inner sep=0pt,
     execute at begin node={\hyphenpenalty=10000\exhyphenpenalty=10000\relax}},
  q/.style={font=\footnotesize\itshape},             % qualifier: multiplicity, mode, arrow label
  sm/.style={font=\footnotesize},                    % small roman label
  blob/.style={draw, rounded corners=7pt, line width=0.6pt},          % stage container
  btitle/.style={fill=white, inner xsep=2pt, inner ysep=1pt, anchor=west}, % title sitting on the blob border
  slab/.style={draw, rounded corners=2.5pt, minimum height=10.5pt, inner xsep=3pt, align=left, anchor=north west},
  ghost/.style={slab, dash pattern=on 1.4pt off 1.1pt},               % "same as Stage 1"
  pill/.style={draw, rounded corners=4.5pt, minimum height=10.5pt, inner xsep=3.5pt, align=center},  % formula pill
  tool/.style={draw, rounded rectangle, minimum height=17pt, inner xsep=5pt, align=center, line width=0.6pt},  % solver / engine
  chev/.style={draw, signal, signal to=east, signal from=west, signal pointer angle=125,
               minimum height=19pt, align=center, inner xsep=1pt, font=\footnotesize},
  chevL/.style={draw, signal, signal to=west, signal from=east, signal pointer angle=125,
               minimum height=19pt, align=center, inner xsep=1pt, font=\footnotesize},
  chevD/.style={draw, signal, signal to=south, signal from=north, signal pointer angle=140,
               minimum width=15pt, minimum height=10pt, align=center, inner sep=0pt, font=\footnotesize\bfseries},
  dspec/.style={fill=\defshade},
  dot/.style={circle, fill, inner sep=0.75pt},
  station/.style={circle, draw, fill=white, minimum size=8.5pt, inner sep=0pt, font=\footnotesize\bfseries},
  box/.style={draw, rounded corners=3pt, align=center, inner xsep=3pt, inner ysep=2.2pt},
}

% ---------- glyphs drawn inside a card centered at the scope origin (card 11 x 14) ----------
\newcommand{\gGraph}{% execution graph: 3 instruction nodes in spo + one rf arc (matches fig a)
  \draw[line width=0.3pt] (-1.8,3.6) -- (-1.8,-3.6);
  \fill (-1.8,3.6) circle (0.95) (-1.8,0) circle (0.95) (-1.8,-3.6) circle (0.95);
  \draw[line width=0.35pt, -{Latex[length=1.5pt,width=1.5pt]}] (-1.0,3.8) .. controls (3.6,3.2) and (3.6,-3.2) .. (-0.9,-3.5);}
\newcommand{\gAsm}{% program text
  \foreach \y/\l in {4/5.2,2/3.8,0/5.2,-2/2.8,-4/4.4} \draw[line width=0.35pt] (-3,\y) -- ++(\l,0);}
\newcommand{\gAnn}{% annotation: a pair of curly braces, drawn (no text)
  \draw[line width=0.4pt, line join=round] (-1.2,3.8) .. controls (-2.4,3.8) and (-2.0,1.0) .. (-2.4,0.3) -- (-3.1,0) -- (-2.4,-0.3) .. controls (-2.0,-1.0) and (-2.4,-3.8) .. (-1.2,-3.8);
  \draw[line width=0.4pt, line join=round] (1.2,3.8) .. controls (2.4,3.8) and (2.0,1.0) .. (2.4,0.3) -- (3.1,0) -- (2.4,-0.3) .. controls (2.0,-1.0) and (2.4,-3.8) .. (1.2,-3.8);}
\newcommand{\gCode}{% source code: </> drawn (no text)
  \draw[line width=0.45pt, line join=miter] (-1.9,0.6) -- (-4.0,-1.2) -- (-1.9,-3.0) (1.9,0.6) -- (4.0,-1.2) -- (1.9,-3.0) (0.9,1.0) -- (-0.9,-3.4);}
\newcommand{\gTrace}{% pipeline trace: staggered stage bars (O3PipeView-like)
  \foreach \y/\a/\b in {4/-3.6/0.4,1.6/-2.4/1.6,-0.8/-1.2/2.8,-3.2/0/3.6} \draw[line width=0.9pt] (\a,\y) -- (\b,\y);}

% ---------- card stack: #1 name, #2 (x,y) of the front card's center, #3 number of cards, #4 glyph macro ----------
% back cards sit up-right of the front card; anchors #1-n (top of the whole stack), #1-s, #1-e, #1-w
\def\cw{11}\def\ch{14}\def\co{1.8}
\newcommand{\cardstack}[4]{%
  \begin{scope}[shift={#2}]
    \foreach \k in {#3,...,1} {\pgfmathsetmacro\o{\co*(\k-1)}
      \draw[fill=white, rounded corners=0.9pt] (-\cw/2+\o,-\ch/2+\o) rectangle ++(\cw,\ch);}
    #4
    \pgfmathsetmacro\top{\ch/2+\co*(#3-1)}\pgfmathsetmacro\rgt{\cw/2+\co*(#3-1)}
    \path (\co*#3/2-\co/2,\top) coordinate (#1-n) (0,-\ch/2) coordinate (#1-s)
          (\rgt,\co*#3/2-\co/2) coordinate (#1-e) (-\cw/2,0) coordinate (#1-w);
  \end{scope}}
% stack of program + annotation PAIRS: each card is two half-cards clipped together
\newcommand{\pairstack}[3]{% #1 name, #2 (x,y), #3 n
  \begin{scope}[shift={#2}]
    \foreach \k in {#3,...,1} {\pgfmathsetmacro\o{\co*(\k-1)}
      \draw[fill=white, rounded corners=0.9pt] (-11+\o,-7+\o) rectangle ++(10.4,14);
      \draw[fill=white, rounded corners=0.9pt] (0.6+\o,-7+\o) rectangle ++(10.4,14);
      \draw[line width=0.5pt] (-1.4+\o,4.5+\o) -- (1.4+\o,4.5+\o);}  % clip joining the two halves
    \begin{scope}[shift={(-5.8,0)}]\gAsm\end{scope}
    \begin{scope}[shift={(5.8,0)}]\gAnn\end{scope}
    \pgfmathsetmacro\top{7+\co*(#3-1)}
    \path (-5.8,7) coordinate (#1-asm-n) (5.8,7) coordinate (#1-ann-n) (-5.8,-7) coordinate (#1-asm-s) (5.8,-7) coordinate (#1-ann-s)
          (11,0) coordinate (#1-e) (-11,0) coordinate (#1-w) (0,\top) coordinate (#1-n) (0,-7) coordinate (#1-s);
  \end{scope}}

% ---------- defense-specific source of truth: gray folded document; #1 name, #2 (x,y), #3 glyph or empty ----------
\newcommand{\sot}[3]{%
  \begin{scope}[shift={#2}]
    \draw[fill=\defshade] (-5.5,-7) -- (5.5,-7) -- (5.5,4) -- (2.5,7) -- (-5.5,7) -- cycle;
    \draw (2.5,7) -- (2.5,4) -- (5.5,4);
    \if\relax\detokenize{#3}\relax
      \foreach \y in {1.8,-0.8,-3.4} \draw[line width=0.3pt] (-3.2,\y) -- (3.2,\y);
    \else #3\fi
    \path (0,7) coordinate (#1-n) (0,-7) coordinate (#1-s) (5.5,0) coordinate (#1-e) (-5.5,0) coordinate (#1-w);
  \end{scope}}

% ---------- bug glyph (#1 position, #2 scale, #3 color) ----------
\newcommand{\buggl}[3]{%
  \begin{scope}[shift={#1}, scale=#2, line width=0.45pt, draw=#3, fill=#3]
    \foreach \y in {1.0,0,-1.0} {\draw (-2.6,\y-0.45) -- (2.6,\y+0.45); \draw (-2.6,\y+0.45) -- (2.6,\y-0.45);}
    \fill (0,-0.2) ellipse (1.35 and 1.9);
    \fill (0,2.15) circle (0.8);
    \draw (-0.35,2.7) -- (-1.2,3.7) (0.35,2.7) -- (1.2,3.7);
  \end{scope}}
% magnifier (#1 center of lens, #2 radius)
\newcommand{\magn}[2]{\draw[line width=0.6pt] #1 circle (#2); \draw[line width=1.1pt] ($#1+(-45:#2)$) -- ++(-45:#2*0.9);}
% legend swatch
\newcommand{\swatch}[1]{\draw[fill=\defshade] #1 rectangle ++(6,6);}

\begin{document}
\begin{tikzpicture}[x=1pt,y=1pt]
% ---------------- knobs ----------------
\def\LX{26}      % centre of the sources-of-truth lane
\def\PLa{54}     % Stage 1 / Stage 3 left edge
\def\PRa{254}    % Stage 1 right edge
\def\PRc{301}    % Stage 3 right edge
\def\PLb{325}    % Stage 1.5 / Stage 2 left edge
\def\PRb{515.5}  % right edge
\def\TOP{-7}     % top of row 1 blobs
\def\BP{4}       % padding below the last element of a blob
\def\XA{318}     % gutter: programs line
\def\XD{310}     % gutter: annotations line
\pgfmathsetmacro\ILa{\PLa+6}\pgfmathsetmacro\ILb{\PLb+6}
\pgfmathsetmacro\CBa{\PRa-16}\pgfmathsetmacro\CBb{\PRb-9}
\pgfmathsetmacro\SLWa{\CBa-5-\ILa-6}   % slab text widths
\pgfmathsetmacro\SLWsa{\SLWa-8+1}
\pgfmathsetmacro\SLWb{\CBb-5-\ILb-6}
\tikzset{
  ftool/.style={draw, rounded rectangle, minimum height=21pt, inner xsep=4pt, align=center, line width=0.6pt, anchor=west},
  passtxt/.style={anchor=west, align=left},
  lang/.style={font=\footnotesize\itshape, anchor=north west},
  sotl/.style={sm, anchor=north, align=center, text width=50pt},
  capt/.style={anchor=north west, align=left},
}
\newcommand{\gety}[2]{\path #1; \pgfgetlastxy{\tmpx}{\tmpy}\pgfmathsetmacro#2{\tmpy}}
\coordinate (lxc) at (\LX,0);
\path (0,0) -- (516,0); % fixes the width at exactly 516pt (figure*)

% ---------------- lane header ----------------
\node[q, anchor=north, align=center, text width=48pt] at (\LX,-1) {Defense\\sources of truth};

% ================= Stage 1 =================
\def\T{\TOP}
\node[lang] at (\ILa+2,\T-4.5) {(Alloy relational modeling language): graph constraints};
\node[slab, text width=\SLWa pt] (a1) at (\ILa,\T-16) {\textbf{1.}~Defense-agnostic execution well-formedness};
\node[slab, text width=\SLWa pt] (a2) at ($(a1.south west)+(0,-2)$) {\textbf{2.}~Defense-agnostic minimality requirements (only load-bearing components specified)};
\node[slab, dspec, text width=\SLWa pt] (a3) at ($(a2.south west)+(0,-2)$) {\textbf{3.}~Defense-specific leakage predicate:};
\begin{scope}[every node/.append style={inner xsep=2.5pt}]
\node[slab, dspec, text width=\SLWsa pt] (d1) at ($(a3.south west)+(8,-2)$) {\textbf{3a.}~Speculation contract: what is speculative};
\node[slab, dspec, text width=\SLWsa pt] (d2) at ($(d1.south west)+(0,-2)$) {\textbf{3b.}~Protection set: protected state + propagation};
\node[slab, dspec, text width=\SLWsa pt] (d3) at ($(d2.south west)+(0,-2)$) {\textbf{3c.}~Leakage function: which operands transmit};
\end{scope}
\coordinate (tx) at ($(a3.south west)+(4,0)$); \draw (tx) |- (d3.west); \draw (tx |- d1.west) -- (d1.west); \draw (tx |- d2.west) -- (d2.west);
\pgfmathsetmacro\TWa{\PRa-5-\ILa}
\node[ftool, minimum width=\TWa pt] (A1) at ($(d3.south -| \ILa,0)+(0,-19)$) {Alloy Analyzer\\[-1pt](Kodkod model finder + SAT backend) + bound};
\draw[line width=0.6pt] (a1.north -| \CBa,0) -- (d3.south -| \CBa,0);
\draw[->] (d3.south -| \CBa,0) -- (A1.north -| \CBa,0);
\coordinate (b1) at ($(A1.south)+(0,-\BP)$);
\draw[blob] (\PLa,\T) rectangle (b1 -| \PRa,0);
\node[btitle] at (\ILa+2,\T) {\textbf{Stage 1}\enspace Abstract test generation};
\sot{ds}{(a3.west -| lxc)}{}
\node[sotl] (dsl) at ($(ds-s)+(0,-1.5)$) {Defense specification\\\textit{manual translation}};
\draw[->] (ds-e) -- (a3.west);

% ================= Stage 1.5 =================
\node[lang] at (\ILb+2,\T-4.5) {(Alloy relational modeling language): constraints};
\node[ghost, text width=\SLWb pt] (bb1) at (\ILb,\T-16) {\textbf{1.}};
\node[q, anchor=east] at ($(bb1.east)+(-3,0)$) {Same as Stage 1};
\node[slab, text width=\SLWb pt] (b2) at ($(bb1.south west)+(0,-2)$) {\textbf{2\kern-0.2pt\raisebox{0.7ex}{$\prime$}.}~Loosened minimality requirements $\Rightarrow$ more flexible execution graphs};
\node[ghost, dspec, text width=\SLWb pt] (bb3) at ($(b2.south west)+(0,-2)$) {\textbf{3.}};
\node[q, anchor=east] at ($(bb3.east)+(-3,0)$) {Same as Stage 1};
\node[slab, text width=\SLWb pt] (b4) at ($(bb3.south west)+(0,-2)$) {\textbf{+\,4.}~Pin: graph contains target Stage-1 graph};
\node[ftool] (A2) at ($(b4.south west)+(0,-19)$) {Alloy Analyzer\\[-1pt](Kodkod model finder + SAT backend) + bound};
\draw[line width=0.6pt] (bb1.north -| \CBb,0) -- (b4.south -| \CBb,0);
\draw[->] (b4.south -| \CBb,0) |- (A2.east);
\coordinate (b2b) at ($(A2.south)+(0,-\BP)$);
\draw[blob] (\PLb,\T) rectangle (b2b -| \PRb,0);
\node[btitle] at (\ILb+2,\T) {\textbf{Stage 1.5}\enspace Abstract test enumeration};

% ---- output 1: across the gap, into the Pin slab ----
\gety{(A1.east)}{\YA}\gety{(b4.west)}{\YP}
\pgfmathsetmacro\GX{\PLb-8}
\draw[->] (A1.east) -- (\GX,\YA) -- (\GX,\YP) -- (b4.west);
\cardstack{o1}{(\PRa+17,\YA)}{3}{\gGraph}
\node[capt, anchor=south west, text width=56pt,
      execute at begin node={\hyphenpenalty=10000\exhyphenpenalty=0\relax}] (c1) at (\PRa+5,\YA+14) {Bounded-exhaustive set of defense-specific leaking execution graphs};

% ---- output 2: down from Stage 1.5 ----
\pgfmathsetmacro\SXb{\ILb+14}
\gety{(b2b)}{\G}
\draw[->] (A2.south -| \SXb,0) -- (\SXb,\G-8.6);
\cardstack{o2}{(\SXb-2.7,\G-19.4)}{4}{\gGraph}
\node[capt, text width=\PRb-\SXb-20 pt] (c2) at (\SXb+14,\G-7.5) {More flexible defense-specific leaking execution graphs\\[0pt]{\itshape Non-exhaustive; multiple per Stage-1 graph; parallelizable}};

% ================= Stage 2 =================
\gety{(c2.south)}{\T}\pgfmathsetmacro\T{\T-9}
\node[lang] at (\SXb+10,\T-4.5) {(Python, LLVM): passes, run per graph};
\coordinate (p0) at (\ILb+7.5,\T-17);
\foreach \i/\t [count=\k from 0] in {
   1/{Order instructions ($\rightarrow$ PCs); collect state},
   2/{Concretize instructions},
   3/{Concretize speculation primitives (branches, STL)},
   4/{Interleave instruction chains},
   5/{Assign registers and memory; reserve registers},
   6/{SSA + LLVM emission with PC labels},
   7/{Annotations; LLVM $\rightarrow$ asm; PCs $\rightarrow$ addresses}} {
  \path ($(p0)+(0,{-(\i-1)*13})$) coordinate (n\i-top);
  \draw[fill=white, line width=0.6pt, line join=miter] ($(n\i-top)+(-7,0)$) -- ($(n\i-top)+(0,-3.6)$) -- ($(n\i-top)+(7,0)$) -- ($(n\i-top)+(7,-10.8)$) -- ($(n\i-top)+(0,-14.4)$) -- ($(n\i-top)+(-7,-10.8)$) -- cycle;
  \coordinate (n\i-c) at ($(n\i-top)+(0,-7.7)$);
  \coordinate (n\i-bot) at ($(n\i-top)+(0,-14.4)$);
  \node[font=\footnotesize\bfseries] at (n\i-c) {\i};
  \node[passtxt] (p\i) at ($(n\i-c)+(10,0)$) {\t};
}
\draw[->] (o2-s) |- ($(n1-top)+(0,3)$) -- ($(n1-top)+(0,-3.6)$);
\coordinate (b3) at ($(n7-bot)+(0,-\BP)$);
\draw[blob] (\PLb,\T) rectangle (b3 -| \PRb,0);
\node[btitle] at (\SXb+8,\T) {\textbf{Stage 2}\enspace Concretization};

% ---- output 3: under Stage 2 ----
\gety{(b3)}{\G}
\gety{(n7-bot)}{\NB}
\draw[->] (n7-bot) -- ($(n7-bot)+(0,-12)$);
\pairstack{o3}{($(n7-bot)+(1.8,-19)$)}{4}
\node[capt, text width=152pt] (c3) at ($(o3-e)+(10,6.5)$) {Programs (x86 asm) + annotations\\[0pt]{\itshape One pair per branch mode\\$\times$ interleave variant}};

% ================= Stage 3 =================
\gety{(b1)}{\T}\pgfmathsetmacro\T{\T-10}
\node[lang] at (\ILa+2,\T-4.5) {(gem5)};
\pgfmathsetmacro\YAs{\T-15}\pgfmathsetmacro\YDs{\T-23}  % programs / directives lines across the top of Stage 3
\pgfmathsetmacro\CT{\T-50}                             % gem5 container top
% gem5 container: logo + "processor model"; harness ring around the gray die
\draw[rounded corners=4pt, line width=0.6pt] (\ILa,\CT) rectangle (284,\CT-73) coordinate (gemSE);
\node[q, anchor=north east] (pm) at (280,\CT-3.5) {processor model};
\node[anchor=east] at ($(pm.west)+(-1.5,0.3)$) {\includegraphics[height=7.5pt]{gem5logo.pdf}};
\draw[rounded corners=3pt] (\ILa+6,\CT-18) rectangle (278,\CT-68);
\node[draw, rounded corners=2pt, dspec, minimum width=272-\ILa-12 pt, minimum height=18pt, anchor=north west] (die) at (\ILa+12,\CT-23) {Defense implementation under test};
\foreach \x in {170,210,250} \draw[line width=0.5pt] (\x,\CT-18) -- (\x,\CT-23);
\node[sm, anchor=north west, align=left, text width=198pt] (hl) at (\ILa+12,\CT-45) {+\,Non-invasive harness: pipeline-stage probes; speculation and timing directives};
% inputs: programs -> die (top), annotations -> ring (speculation directives) and -> Inspector (labels)
\gety{(o3-w)}{\YO}
\draw[->] (o3-w) -- (\XA,\YO) -- (\XA,\YAs) -- (100,\YAs) -- (100,0 |- die.north);
\gety{(o3-ann-s)}{\YN}\pgfmathsetmacro\YJ{\YN-5}
\draw[->] (o3-ann-s) -- (o3-ann-s |- 0,\YJ) -- (\XD,\YJ) -- (\XD,\YDs) -- (128,\YDs) -- (128,\CT-18);
\node[q, anchor=north west] at (131,\YDs-2) {speculation directives};
\node[dot] at (\XD,\YJ+12) {};
% pipeline info artifact -> inspector -> bug
\pgfmathsetmacro\YI{\YJ+12}
\cardstack{tr}{(195,\YI)}{1}{\gTrace}
\draw[->] (195,\CT-73) -- (tr-n);
\node[sm, anchor=east, align=right] at ($(tr-w)+(-3,0)$) {Pipeline info\\artifact};
\node[box, anchor=west, minimum height=15pt, align=left, inner xsep=3pt] (I) at (215,\YI) {\hspace*{8pt}Inspector\hspace*{3pt}};
\magn{($(I.west)+(6,0.8)$)}{2.5}
\draw[->] (tr-e) -- (tr-e -| I.west);
\draw[->] (\XD,\YI) -- (I.east) node[q, pos=0.45, anchor=south, yshift=1.5pt] {labels};
\pgfmathsetmacro\YB{\YI-21}
\draw[->, line width=0.7pt] (I.south -| 253,0) -- (253,\YB+6.5);
\buggl{(253,\YB)}{1.6}{black}
\node[font=\footnotesize\bfseries, anchor=west] at (261.5,\YB) {Bug?};
% sources: defense implementation -> die ; transmission mechanisms -> harness ("selects hooks") + inspector
\sot{gb}{(die.west -| lxc)}{\gCode}
\node[sotl] (gbl) at ($(gb-s)+(0,-1.5)$) {Defense implementation (gem5)};
\draw[->] (gb-e) -- (die.west);
\draw[dashed, ->] ($(dsl.south)+(0,-2)$) -- ($(gb-n)+(0,2)$) coordinate[pos=0.5] (gm);
\node[q, rotate=90, anchor=south, align=center] at ($(gm)+(-2,0)$) {\textbf{SimSpect goal:} validate security\\guarantee on implementation};
\sot{tm}{(\LX,\YB+25)}{}
\node[sotl] (tml) at ($(tm-s)+(0,-1.5)$) {Defense transmission mechanisms};
\draw (tm-e) -- (\PLa+4,0 |- tm-e) |- (229,\YB+1);
\draw[->] (229,\YB+1) -- (229,0 |- I.south);
\draw[->] (96,\YB+1) -- (96,\CT-68) node[q, pos=0.45, anchor=east, align=right, xshift=-2pt] {selects\\hooks};
\node[dot] at (96,\YB+1) {};
\draw[blob] (\PLa,\T) rectangle (\PRc,\YB-9);
\node[btitle] at (\ILa+2,\T) {\textbf{Stage 3}\enspace Execution};
\end{tikzpicture}
\end{document}
